% !TEX root = ../main.tex
\corrigesChapitre{1}{Composition d'un système chimique}
\corrigesSouples

\begin{blocLivre}{cor-1-1}
\begin{corrige}{1}{Calculer des masses molaires}
\demarche{La masse molaire d'une espèce est la somme des masses molaires atomiques de tous les atomes de l'entité.}
\begin{questions}
\item $M(\ce{C8H10N4O2})=8\times12{,}0+10\times1{,}0+4\times14{,}0+2\times16{,}0=\res{\qty{194,0}{g.mol^{-1}}}$ ;
$M(\ce{C6H12O6})=6\times12{,}0+12\times1{,}0+6\times16{,}0=\res{\qty{180,0}{g.mol^{-1}}}$.
\item $M(\ce{C9H8O4})=9\times12{,}0+8\times1{,}0+4\times16{,}0=\res{\qty{180,0}{g.mol^{-1}}}$ : la même valeur que le glucose. Deux espèces de formules
différentes peuvent avoir la \res{\text{même masse molaire}} : $M$ ne suffit pas à identifier une espèce.
\item $M(\ce{SO4^{2-}})=32{,}1+4\times16{,}0=\res{\qty{96,1}{g.mol^{-1}}}$. La masse d'un électron est environ \num{2000} fois plus petite que celle d'un
nucléon : elle est \res{\text{négligeable}}.
\item $M(\ce{NH3})=14{,}0+3\times1{,}0=\res{\qty{17,0}{g.mol^{-1}}}$.
\end{questions}
\end{corrige}
\end{blocLivre}
\begin{blocLivre}{cor-1-2}
\begin{corrige}{2}{Un diamant et une bague}
\demarche{$n=\dfrac{m}{M}$, puis $N=n\times N_\mathrm{A}$.}
\begin{questions}
\item $n(\ce{C})=\dfrac{0{,}200}{12{,}0}=\res{\qty{1,67e-2}{mol}}$ et $N=\num{1.67e-2}\times\num{6.02e23}=\res{\num{1.00e22}}$ atomes.
\item $n(\ce{Au})=\dfrac{5{,}0}{197{,}0}=\res{\qty{2,54e-2}{mol}}$ et $N=\res{\num{1.53e22}}$ atomes.
\item La \res{\text{bague}} contient le plus d'atomes, mais seulement 1,5 fois plus pour une masse 25 fois plus grande : un atome
d'or est bien plus lourd qu'un atome de carbone. Pour comparer des nombres d'entités, il faut comparer des quantités de
matière, pas des masses.
\item $m(\text{atome})=\dfrac{M}{N_\mathrm{A}}=\dfrac{12{,}0}{\num{6.02e23}}=\res{\qty{1,99e-23}{g}}$.
\end{questions}
\retiens{Pour comparer des nombres d'entités, on compare des quantités de matière, jamais des masses.}
\end{corrige}
\end{blocLivre}
\begin{blocLivre}{cor-1-3}
\begin{corrige}{3}{Un liquide : l'éthanol}
\demarche{Pour un liquide, $m=\rho\times V$, puis $n=\dfrac{m}{M}$.}
\begin{questions}
\item $m=\rho\times V=0{,}789\times25{,}0=\res{\qty{19,7}{g}}$.
\item $M(\ce{C2H6O})=2\times12{,}0+6\times1{,}0+16{,}0=\qty{46,0}{g.mol^{-1}}$ et $n=\dfrac{19{,}7}{46{,}0}=\res{\qty{0,429}{mol}}$.
\item $N=0{,}429\times\num{6.02e23}=\res{\num{2.58e23}}$ molécules.
\item Une mole a une masse de \qty{46,0}{g}, soit un volume $V=\dfrac{46{,}0}{0{,}789}=\res{\qty{58,3}{mL}}$.
\end{questions}
\end{corrige}
\end{blocLivre}
\begin{blocLivre}{cor-1-4}
\begin{corrige}{4}{Deux ballons de baudruche}
\demarche{$n=\dfrac{V}{V_\mathrm{m}}$ pour un gaz, quel qu'il soit ; puis $m=n\times M$.}
\begin{questions}
\item $n=\dfrac{4{,}8}{24{,}0}=\res{\qty{0,20}{mol}}$ dans chaque ballon : à même température et même pression, un même volume de gaz contient
la même quantité de matière, \res{\text{quelle que soit la nature du gaz}}.
\item $m(\ce{He})=0{,}20\times4{,}0=\res{\qty{0,80}{g}}$ et $m(\ce{CO2})=0{,}20\times44{,}0=\res{\qty{8,8}{g}}$.
\item $m(\text{air})=0{,}20\times29{,}0=\res{\qty{5,8}{g}}$.
\item L'hélium est plus léger que l'air qu'il remplace, le dioxyde de carbone plus lourd : le ballon d'hélium s'élève, l'autre
tombe (la masse de l'enveloppe joue aussi un rôle).
\end{questions}
\schema[corrige]{01-volumemolaire}
\end{corrige}
\end{blocLivre}
\begin{blocLivre}{cor-1-5}
\begin{corrige}{5}{Le sérum physiologique}
\demarche{$t=\dfrac{m}{V}$ et $c=\dfrac{t}{M}$ ; dans un volume donné, $m=t\times V$ et $n=c\times V$.}
\begin{questions}
\item $t=\dfrac{0{,}9}{0{,}100}=\res{\qty{9,0}{g.L^{-1}}}$.
\item $c=\dfrac{t}{M}=\dfrac{9{,}0}{58{,}5}=\res{\qty{0,154}{mol.L^{-1}}}$ : on divise une masse par litre par une masse par mole, et l'on obtient des moles
par litre.
\item $m=9{,}0\times0{,}500=\res{\qty{4,5}{g}}$ et $n=\dfrac{4{,}5}{58{,}5}=\res{\qty{7,69e-2}{mol}}$.
\item Chaque unité $\ce{NaCl}$ libère un ion $\ce{Na+}$ : $N=\num{7.69e-2}\times\num{6.02e23}=\res{\num{4.63e22}}$ ions sodium.
\end{questions}
\piege{Confondre la concentration en masse $t$ d'un soluté et la masse volumique $\rho$ de la solution : toutes deux s'expriment en
\unit{g.L^{-1}}, mais $t$ ne compte que le soluté, alors que $\rho$ compte toute la solution (environ \qty{1000}{g.L^{-1}} ici).}
\end{corrige}
\end{blocLivre}
\begin{blocLivre}{cor-1-6}
\begin{corrige}{6}{Préparer une solution par dissolution}
\demarche{$n=c\times V$, puis $m=n\times M$ ; protocole : pesée, fiole jaugée, dissolution, trait de jauge, homogénéisation.}
\begin{questions}
\item $M=63{,}5+32{,}1+4\times16{,}0+5\times(2\times1{,}0+16{,}0)=\res{\qty{249,6}{g.mol^{-1}}}$.
\item $n=c_0\times V=0{,}0400\times0{,}2500=\res{\qty{1,00e-2}{mol}}$, soit $m=n\times M=\num{1.00e-2}\times249{,}6=\res{\qty{2,50}{g}}$.
\item Peser \qty{2,50}{g} de cristaux dans une coupelle sur une balance tarée ; les introduire dans une \res{\text{fiole jaugée de 250,0 mL}}
avec un entonnoir ; rincer coupelle et entonnoir à l'eau distillée ; remplir aux trois quarts et agiter jusqu'à
dissolution ; compléter jusqu'au trait de jauge, boucher et homogénéiser.
\item $t=c_0\times M=0{,}0400\times249{,}6=\res{\qty{9,98}{g.L^{-1}}}$.
\end{questions}
\schema[corrige]{01-dissolution}
\end{corrige}
\end{blocLivre}
\begin{blocLivre}{cor-1-7}
\begin{corrige}{7}{Diluer une solution}
\demarche{Lors d'une dilution, la quantité de soluté se conserve : $c_0\,V_0=c_1\,V_1$ et $F=\dfrac{c_0}{c_1}=\dfrac{V_1}{V_0}$.}
\begin{questions}
\item $F=\dfrac{0{,}0400}{\num{2.00e-3}}=\res{20}$.
\item $V_0=\dfrac{V_1}{F}=\dfrac{100{,}0}{20}=\res{\qty{5,00}{mL}}$.
\item La \res{\text{pipette jaugée de 5,00 mL}} (C) pour prélever S$_0$, versée dans la \res{\text{fiole jaugée de 100,0 mL}} (E), et un bécher (A) pour
contenir un peu de S$_0$. L'éprouvette n'est pas assez précise, et la fiole de \qty{50,0}{mL} ne donne pas le bon volume.
\item S$_1$ est 20 fois moins concentrée, donc \res{\text{moins colorée}} : $t_1=c_1\times M=\num{2.00e-3}\times249{,}6=\res{\qty{0,499}{g.L^{-1}}}$.
\end{questions}
\schema[corrige]{01-ex07}
\piege{Prélever la solution mère à l'éprouvette graduée : sa précision est insuffisante pour une dilution. On prélève
toujours à la pipette jaugée et l'on complète dans une fiole jaugée.}
\end{corrige}
\end{blocLivre}
\begin{blocLivre}{cor-1-8}
\begin{corrige}{8}{Le degré alcoolique}
\demarche{Volume d'éthanol, puis $m=\rho\times V$, puis $n=\dfrac{m}{M}$.}
\begin{questions}
\item Un litre de vin contient $12{,}5\times10=\qty{125}{mL}$ d'éthanol, soit $m=0{,}789\times125=\res{\qty{98,6}{g}}$.
\item $n=\dfrac{98{,}6}{46{,}0}=\qty{2,14}{mol}$ dans un litre : $\res{c=\qty{2,14}{mol.L^{-1}}}$.
\item \qty{10}{cL} contiennent \qty{12,5}{mL} d'éthanol, soit $0{,}789\times12{,}5=\res{\qty{9,9}{g}}$.
\item Bière : $0{,}05\times250=\qty{12,5}{mL}$, soit \res{\qty{9,9}{g}} ; whisky : $0{,}40\times30=\qty{12}{mL}$, soit \res{\qty{9,5}{g}}. Les verres servis
habituellement contiennent tous environ \qty{10}{g} d'alcool : c'est le \res{\text{verre standard}}, qui sert à compter sa consommation quelle
que soit la boisson.
\end{questions}
\end{corrige}
\end{blocLivre}
\begin{blocLivre}{cor-1-9}
\begin{corrige}{9}{Une eau minérale riche en magnésium}
\demarche{$c=\dfrac{t}{M}$ avec $t$ en \unit{g.L^{-1}} ; $n=c\times V$, puis $N=n\times N_\mathrm{A}$.}
\begin{questions}
\item $c(\ce{Mg^{2+}})=\dfrac{0{,}119}{24{,}3}=\res{\qty{4,90e-3}{mol.L^{-1}}}$ et $c(\ce{Ca^{2+}})=\dfrac{0{,}549}{40{,}1}=\res{\qty{1,37e-2}{mol.L^{-1}}}$.
\item $n=\num{4.90e-3}\times0{,}20=\qty{9,79e-4}{mol}$, soit $N=\res{\num{5.90e20}}$ ions magnésium.
\item $V=\dfrac{300}{119}=\res{\qty{2,5}{L}}$ : c'est beaucoup, l'alimentation apporte heureusement l'essentiel du magnésium.
\item Les apports nutritionnels sont exprimés en masse (milligrammes par jour) : la concentration en masse permet de les
comparer directement, sans connaître la notion de mole.
\end{questions}
\end{corrige}
\end{blocLivre}
\begin{blocLivre}{cor-1-10}
\begin{corrige}{10}{L'air d'une salle de classe}
\demarche{Volume de la salle, puis $n=\dfrac{V}{V_\mathrm{m}}$ ; pour un gaz, les pourcentages en volume sont aussi des pourcentages en quantité de matière.}
\begin{questions}
\item $V=8{,}0\times7{,}0\times3{,}0=\qty{168}{m^3}=\qty{1,68e5}{L}$, d'où $n=\dfrac{\num{1.68e5}}{24{,}0}=\res{\qty{7,00e3}{mol}}$ d'air.
\item $n(\ce{O2})=0{,}21\times\num{7.00e3}=\qty{1,47e3}{mol}$, soit $m=\num{1.47e3}\times32{,}0=\qty{4,70e4}{g}=\res{\qty{47,0}{kg}}$ ;
$n(\ce{N2})=0{,}78\times\num{7.00e3}=\qty{5,46e3}{mol}$, soit \res{\qty{153}{kg}}.
\item $V=30\times12=\qty{360}{L}$, soit $n=\dfrac{360}{24{,}0}=\res{\qty{15}{mol}}$ de dioxyde de carbone.
\item $\dfrac{360}{\num{1.68e5}}=\num{2.1e-3}$, soit \res{\qty{0,21}{\percent}} ou \res{\num{2.1e3}~\text{ppm}}. L'air extérieur en contient environ \qty{420}{ppm} : sans
aération, la teneur dépasse vite les valeurs recommandées et l'air devient inconfortable, d'où la nécessité d'aérer.
\end{questions}
\retiens{Pour un mélange de gaz, un pourcentage en volume est aussi un pourcentage en quantité de matière, car tous les gaz
ont le même volume molaire dans les mêmes conditions.}
\end{corrige}
\end{blocLivre}
\begin{blocLivre}{cor-1-11}
\begin{corrige}{11}{Un programme pour les masses molaires}
\demarche{Un dictionnaire associe une clé (le symbole de l'élément) à une valeur ; la boucle parcourt les clés de \code{formule}.}
\begin{questions}
\item \code{M\_atome} associe à chaque élément sa masse molaire atomique, \code{formule} le nombre d'atomes de chaque élément dans
la molécule. La boucle ajoute, pour chaque élément, le produit du nombre d'atomes par la masse molaire atomique : elle
calcule la \res{\text{masse molaire}} de la molécule.
\item $M=\qty{194,0}{g.mol^{-1}}$, $\res{n=\qty{4,124e-4}{mol}}$ de caféine dans la tasse et $\res{N=\num{2.482e20}}$ molécules.
\item Il suffit de remplacer la ligne \code{formule = \{'C': 6, 'H': 12, 'O': 6\}} (et la masse \code{m}) : le programme affiche
$\res{M=\qty{180,0}{g.mol^{-1}}}$.
\item Il a oublié l'atome d'azote, \code{'N': 1}. Le programme donnerait \qty{137,0}{g.mol^{-1}} au lieu de $\res{\qty{151,0}{g.mol^{-1}}}$.
\end{questions}
\end{corrige}
\end{blocLivre}
\begin{blocLivre}{cor-1-12}
\begin{corrige}{12}{La glycémie}
\demarche{$t=c\times M$ et $c=\dfrac{t}{M}$, en convertissant les millimoles en moles.}
\begin{questions}
\item $t=\num{7.8e-3}\times180{,}0=\res{\qty{1,40}{g.L^{-1}}}$.
\item $\dfrac{0{,}70}{180{,}0}=\qty{3,9e-3}{mol.L^{-1}}$ et $\dfrac{1{,}10}{180{,}0}=\qty{6,1e-3}{mol.L^{-1}}$ : on retrouve l'intervalle \res{\qtyrange{3,9}{6,1}{mmol.L^{-1}}}.
\item La glycémie dépasse le seuil de \qty{1,26}{g.L^{-1}}, soit \qty{7,0}{mmol.L^{-1}}. Mais le diagnostic exige \res{\text{deux dosages}} : un seul
résultat ne suffit pas. Il faut consulter un médecin, qui prescrira un second dosage à jeun.
\item $m=1{,}40\times5{,}0=\res{\qty{7,0}{g}}$ de glucose, contre $1{,}10\times5{,}0=\qty{5,5}{g}$ à la limite haute : l'excès vaut \res{\qty{1,5}{g}}, environ un quart de
morceau de sucre. Une faible masse suffit à dérégler la glycémie.
\end{questions}
\end{corrige}
\end{blocLivre}
\begin{blocLivre}{cor-1-13}
\begin{corrige}{13}{Le réservoir d'une voiture à hydrogène}
\demarche{$n=\dfrac{m}{M}$ avec $m$ en grammes ; $V=n\times V_\mathrm{m}$ dans les conditions usuelles.}
\begin{questions}
\item $n=\dfrac{\num{5600}}{2{,}0}=\res{\qty{2800}{mol}}$, soit $N=2800\times\num{6.02e23}=\res{\num{1.69e27}}$ molécules.
\item $V=2800\times24{,}0=\qty{67200}{L}$, soit \res{\qty{67}{m^3}}, environ \res{473} fois le volume des réservoirs : sans compression, le gaz
occuperait le volume d'une pièce entière.
\item $n=\dfrac{800}{2{,}0}=\qty{400}{mol}$, soit $V=400\times24{,}0=\res{\qty{9600}{L}}$.
\item $\dfrac{5{,}6}{0{,}8}\times100=\res{\qty{700}{km}}$, un peu plus que les \qty{650}{km} annoncés : tout le dihydrogène n'est pas utilisable (une
réserve reste dans les réservoirs) et la consommation dépend de la conduite.
\end{questions}
\end{corrige}
\end{blocLivre}
\begin{blocLivre}{cor-1-14}
\begin{corrige}{14}{La préparation d'une perfusion}
\demarche{$m=t\times V$, $n=\dfrac{m}{M}$ ; après l'ajout, le volume de la poche vaut \qty{510}{mL}.}
\begin{questions}
\item $m=50\times0{,}500=\res{\qty{25}{g}}$ et $n=\dfrac{25}{180{,}0}=\res{\qty{0,139}{mol}}$.
\item $M(\ce{KCl})=39{,}1+35{,}5=\qty{74,6}{g.mol^{-1}}$, $n=\dfrac{1{,}0}{74{,}6}=\qty{1,34e-2}{mol}$ : \res{\qty{13,4}{mmol}} d'ions potassium.
\item $c(\ce{K+})=\dfrac{\num{1.34e-2}}{0{,}510}=\res{\qty{2,63e-2}{mol.L^{-1}}}$, soit \qty{26,3}{mmol.L^{-1}} ; $t(\text{glucose})=\dfrac{25}{0{,}510}=\res{\qty{49,0}{g.L^{-1}}}$.
\item \qty{26,3}{mmol.L^{-1}} est inférieur à \qty{40}{mmol.L^{-1}} : la consigne est \res{\text{respectée}}. Avec deux ampoules,
$c=\dfrac{2\times13{,}4}{520}=\qty{51,5}{mmol.L^{-1}}$, trop : on ne peut ajouter qu'\res{\text{une seule ampoule}}.
\end{questions}
\end{corrige}
\end{blocLivre}
\begin{blocLivre}{cor-1-15}
\begin{corrige}{15}{De l'or 18 carats}
\demarche{Masse de chaque métal à partir des pourcentages massiques, puis $n=\dfrac{m}{M}$ ; comparer les quantités de matière.}
\begin{questions}
\item \qty{24}{g} d'alliage contiennent $0{,}75\times24=\res{\qty{18}{g}}$ d'or : l'alliage titre bien 18 carats.
\item $m(\ce{Au})=\qty{3,0}{g}$, $m(\ce{Ag})=m(\ce{Cu})=\qty{0,50}{g}$ ; $n(\ce{Au})=\res{\qty{1,52e-2}{mol}}$, $n(\ce{Ag})=\res{\qty{4,63e-3}{mol}}$ et
$n(\ce{Cu})=\res{\qty{7,87e-3}{mol}}$.
\item $n_\text{total}=\qty{2,77e-2}{mol}$ : la part des atomes d'or vaut $\dfrac{1{,}52}{2{,}77}=\res{\qty{55}{\percent}}$. Le client a tort : un peu plus d'un atome
sur deux seulement est un atome d'or.
\item Le \res{\text{cuivre}} (28~\% des atomes) devance l'argent (17~\%), alors que leurs masses sont égales : un atome de cuivre est plus
léger qu'un atome d'argent, donc une même masse en contient davantage.
\end{questions}
\piege{Un pourcentage en masse n'est pas un pourcentage en nombre d'atomes : pour passer de l'un à l'autre, il faut diviser
chaque masse par la masse molaire correspondante.}
\end{corrige}
\end{blocLivre}
